\documentclass[11pt, a4paper, article]{article}

%========================================================================================
% PREAMBLE
%========================================================================================

%----- Page Layout -----
\usepackage[margin=1in, headheight=13.6pt]{geometry}

%----- Fonts and Symbols -----
\usepackage{amsmath}
\usepackage{amssymb}
\usepackage{amsfonts}
\usepackage{bm} % For bold math symbols (e.g., vectors)
\usepackage{graphicx}
\usepackage{lmodern} % For a modern font
\usepackage[T1]{fontenc}
\usepackage{textcomp}

%----- Headers and Footers -----
\usepackage{fancyhdr}

%----- Hyperlinks and Referencing -----
\usepackage{hyperref}
\usepackage{xcolor}

%----- Document Formatting -----
\usepackage{abstract}
\usepackage{titling}
\usepackage{parskip} % Adds space between paragraphs instead of indentation

%========================================================================================
% CUSTOM COMMANDS AND SETTINGS
%========================================================================================

%----- Math Commands -----
\renewcommand{\vec}[1]{\bm{#1}} % Redefine \vec for bold vectors
\newcommand{\pd}[2]{\frac{\partial #1}{\partial #2}} % Partial derivative
\newcommand{\dd}[2]{\frac{d #1}{d #2}} % Total derivative
\newcommand{\funcint}[1]{\mathcal{D}[#1]} % Functional integral measure

%----- Hyperlink Setup -----
\hypersetup{
    colorlinks=true,
    linkcolor=blue!50!black,
    citecolor=green!50!black,
    urlcolor=cyan!50!black,
    pdftitle={QFT Lecture 1: From Classical Particles to Classical Fields},
    pdfauthor={Course: QFT 2025 Part I},
    pdfsubject={Lecture notes for an introductory course in Quantum Field Theory},
    pdfkeywords={QFT, Field Theory, Statistical Mechanics, Phase Transitions, Continuum Limit}
}

%----- Header and Footer Setup -----
\pagestyle{fancy}
\fancyhf{} % Clear all header and footer fields
\fancyhead[L]{QFT 2025 Part I}
\fancyhead[R]{Lecture 1: From Classical Particles to Classical Fields}
\fancyfoot[C]{\thepage}
\renewcommand{\headrulewidth}{0.4pt}
\renewcommand{\footrulewidth}{0.4pt}

%----- Title Setup -----
\pretitle{\begin{center}\Huge\bfseries}
\posttitle{\end{center}}
\preauthor{\begin{center}\large}
\postauthor{\end{center}}
\predate{\begin{center}\large}
\postdate{\end{center}}

%========================================================================================
% DOCUMENT START
%========================================================================================

\title{Lecture 1 From Classical Particles to Classical Fields}
\author{Course: QFT 2025 Part I \\ \small{\url{https://nimalashkari.com/quantum-field-theory-2025-part-i/}}}
\date{\today}

\begin{document}

\maketitle
\thispagestyle{fancy}

\begin{abstract}
\noindent This lecture provides a pedagogical introduction to the core ideas of why it is natural to replace particles with fields as fundamental building blocks of many-body systems, both at the classical and quantum levels. We begin by motivating the transition from a particle-based description of nature to a field-based one. We will argue that field theory is a natural theoretical framework for modeling many-body systems, and it is relevant both at high energies (particle physics) and low energies near a continuous phase transition (condensed matter physics). We will develop this idea using a concrete example from statistical mechanics —the classical XY model of magnetism. By analyzing this system, we will explicitly derive a classical field theory as the effective description of long-wavelength physics. This journey will illuminate the foundational principles common to all field theories: the role of fluctuations, locality, symmetry, and the concept of universality.
\end{abstract}

\newpage
\tableofcontents
\newpage

\section{Motivation: Why Fields?}

\subsection{The Limits of the Particle Picture}
Our classical intuition, shaped by Newtonian mechanics, is built upon the concept of \textbf{particles}. A system, be it a collection of billiard balls or planets, is described by the set of positions and momenta of its constituent parts: $\{ \vec{q}_i(t), \vec{p}_i(t) \}_{i=1}^N$. The state of the system is a single point in a $6N$-dimensional phase space, and its evolution is governed by Hamilton's equations. This framework is deterministic, intuitive, and incredibly successful for systems with a small number of degrees of freedom.

However, this picture breaks down in two important regimes:
\begin{enumerate}
    \item \textbf{Many-Body Systems:} When $N$ becomes large, as in a solid, liquid, or gas ($N \sim 10^{23}$), tracking every particle is both impossible and pointless. We are not interested in the exact trajectory of the $10^{15}$-th molecule in a glass of water. Instead, we care about its collective, macroscopic properties like pressure, density, and temperature. These are inherently continuous quantities, best described by \textbf{fields}. Fluid dynamics, for example, abandons the particle picture entirely in favor of the density field $\rho(\vec{x}, t)$ and the velocity field $\vec{v}(\vec{x}, t)$.

    \item \textbf{Relativity and Quantum Mechanics:} The classical particle is a localized point object. This concept is fundamentally at odds with both special relativity (which forbids instantaneous action at a distance, making field-mediated interactions necessary) and quantum mechanics (where the Heisenberg uncertainty principle implies that a perfectly localized particle would have infinite momentum uncertainty). Quantum mechanics introduced the \textbf{wave function}, $\psi(\vec{x},t)$, which is a field that describes the probability of finding a particle.
\end{enumerate}

\subsection{Quantum Field Theory as a Synthesis}
\textbf{Quantum Field Theory (QFT)} is the framework that reconciles and unifies these ideas. It starts with fields as the fundamental entities of nature. In QFT:
\begin{itemize}
    \item The universe is filled with a set of fundamental fields (the electron field, the photon field, the Higgs field, etc.).
    \item What we perceive as \textbf{particles} are quantized excitations—ripples or vibrations—of these underlying fields.
\end{itemize}
This perspective elegantly resolves the particle-wave duality: the entity is fundamentally a field, but its quantized excitations are discrete and behave like particles. It also naturally incorporates the creation and annihilation of particles, a cornerstone of high-energy physics, as simply the process of adding or removing energy from the field's vibrational modes.

While QFT is the language of the Standard Model of Particle Physics, its principles are far more general. The core idea—that a continuous field description can emerge from an underlying discrete system—is a universal phenomenon. We will now explore this emergence in a more familiar, classical context.

\newpage
\section{A Toy Model: The Classical XY Model}
To build our intuition, we will analyze a concrete system from statistical mechanics: a model of magnetism in a crystalline solid. This will serve as our "laboratory" for seeing a field theory emerge.

\subsection{Defining the Model}
Imagine a $d$-dimensional hypercubic lattice. At each discrete site, indexed by $i$ (with position vector $\vec{x}_i$), sits a tiny magnetic moment. For simplicity, we assume this moment is constrained to rotate in a two-dimensional plane. We can represent it as a unit vector, $\vec{s}_i$, called a "spin" or "rotor".

The orientation of each spin can be described by a single angle $\theta_i \in [0, 2\pi)$:
\begin{equation}
    \vec{s}_i = (\cos\theta_i, \sin\theta_i)
\end{equation}
The physics of the system is determined by the interactions between these spins. The dominant interaction is typically a short-range force where neighboring spins prefer to align. This minimizes the energy. A simple mathematical model for this is the \textbf{XY model}, whose Hamiltonian (energy function) is:
\begin{equation}
    H = -J \sum_{\langle i,j \rangle} \vec{s}_i \cdot \vec{s}_j = -J \sum_{\langle i,j \rangle} \cos(\theta_i - \theta_j)
    \label{eq:hamiltonian}
\end{equation}
Here, $J > 0$ is the ferromagnetic coupling constant that sets the energy scale of the interaction, and the sum $\sum_{\langle i,j \rangle}$ runs over all pairs of nearest-neighbor sites on the lattice. The minus sign ensures that the energy is lowest when the spins are aligned ($\theta_i - \theta_j = 0$).



\subsection{Limiting Cases: Order vs. Disorder}
The state of our system is determined by the competition between the ordering tendency of the interaction energy ($J$) and the disordering tendency of thermal fluctuations ($k_B T$).

\begin{itemize}
    \item \textbf{Zero Temperature ($T=0$):} In this limit, the system will settle into its lowest possible energy state (the ground state). To minimize $H$, we must maximize every $\cos(\theta_i - \theta_j)$ term. This occurs when all spins are perfectly aligned, i.e., $\theta_i = \theta_j$ for all neighbors. This implies that all spins in the entire system point in the same, spontaneously chosen direction. This is a state of perfect \textbf{long-range order}, characteristic of a \textbf{ferromagnet}. The system has a non-zero net magnetization.

    \item \textbf{Infinite Temperature ($T \to \infty$):} In this limit, thermal energy $k_B T$ is far greater than the coupling energy $J$. The Boltzmann factor, $e^{-H/k_B T}$, approaches $e^0 = 1$ for all configurations. This means every possible arrangement of spins is equally likely. The spins will be oriented randomly, and on average, there is no net magnetization. This is a state of complete \textbf{disorder}, characteristic of a \textbf{paramagnet}.
\end{itemize}

Somewhere between these two extremes, for dimensions $d > 2$, a \textbf{phase transition} must occur.

\newpage
\section{Statistical Mechanics and Critical Phenomena}

\subsection{Correlations and Correlation Length}
To describe the system at a finite temperature, we need the tools of statistical mechanics. We are interested in average quantities. A crucial diagnostic is the \textbf{two-point correlation function}:
\begin{equation}
    G(\vec{r}_{ij}) = \langle \vec{s}_i \cdot \vec{s}_j \rangle \equiv \langle \cos(\theta_i - \theta_j) \rangle
\end{equation}
where $\vec{r}_{ij} = \vec{x}_i - \vec{x}_j$ is the separation vector, and the angle brackets $\langle \dots \rangle$ denote a thermal average over all possible configurations, weighted by the Boltzmann factor:
\begin{equation}
    \langle \mathcal{O} \rangle = \frac{1}{Z} \sum_{\text{configs}} \mathcal{O}(\text{config}) e^{-H(\text{config})/k_B T}
\end{equation}
Here, $Z$ is the partition function, which normalizes the probabilities.

The correlation function measures how strongly the orientation of a spin at site $i$ influences the orientation of a spin at site $j$. In a disordered phase (high temperature), we expect this influence to die off quickly with distance. Typically, it decays exponentially:
\begin{equation}
    G(r) \sim e^{-r/\xi} \quad \text{for large } r = |\vec{r}|
\end{equation}
The parameter $\xi$ is the \textbf{correlation length}. It is the characteristic length scale of the ordered fluctuations in the system. Beyond this distance, the spins are effectively decorrelated.

\subsection{The Continuous Phase Transition}
As we lower the temperature from a very high value, the correlation length $\xi$ grows. Spins start to form larger and larger correlated patches. A \textbf{continuous phase transition} (or second-order phase transition) occurs at a critical temperature, $T_c$, where something remarkable happens:
\begin{equation}
    \xi(T) \to \infty \quad \text{as} \quad T \to T_c
\end{equation}
At the critical point, the system is \textbf{scale-invariant}. Correlated fluctuations exist on \textit{all length scales}, from the microscopic lattice spacing $a$ all the way up to the size of the system. The system looks statistically the same at any level of magnification.



This divergence is the key insight. When we are close to the critical point, the correlation length $\xi$ is enormous compared to the microscopic lattice spacing $a$. The physics that we observe at long distances is dominated by the collective behavior of huge clusters of spins. The discrete, microscopic nature of the underlying lattice becomes irrelevant. This is the perfect setting for a continuum description to emerge.

\newpage
\section{Deriving the Continuum Field Theory}

We will now formalize the transition from our discrete spin model to a continuous field theory. The goal is to find an effective Hamiltonian (or, more appropriately, a free energy functional) that describes the long-wavelength physics near $T_c$.

\subsection{The Low-Temperature (Slowly Varying) Approximation}
Near the critical point (and especially in the low-temperature phase), the spins are mostly aligned over large regions. This means that for two adjacent sites $i$ and $j$, the difference in their angles, $\Delta\theta = \theta_i - \theta_j$, is very small. This is our key physical assumption.

Let us define a continuous field $\theta(\vec{x})$ that smoothly interpolates the discrete angles $\theta_i$ on the lattice, such that $\theta(\vec{x}_i) = \theta_i$. Consider two neighboring sites, with site $j$ at position $\vec{x}_j = \vec{x}_i + a\hat{n}$, where $a$ is the lattice spacing and $\hat{n}$ is a unit vector pointing from $i$ to $j$. We can Taylor expand the field $\theta(\vec{x})$ around the point $\vec{x}_i$:
\begin{equation}
    \theta_j = \theta(\vec{x}_i + a\hat{n}) \approx \theta(\vec{x}_i) + a (\hat{n} \cdot \vec{\nabla}\theta(\vec{x}_i)) + \mathcal{O}(a^2)
\end{equation}
The difference in angle is therefore approximately:
\begin{equation}
    \theta_i - \theta_j \approx -a (\hat{n} \cdot \vec{\nabla}\theta(\vec{x}_i))
\end{equation}
Since this difference is small, we can expand the cosine term in the Hamiltonian (\ref{eq:hamiltonian}) using the standard approximation $\cos(x) \approx 1 - x^2/2$:
\begin{equation}
    \cos(\theta_i - \theta_j) \approx 1 - \frac{1}{2}(\theta_i - \theta_j)^2 \approx 1 - \frac{a^2}{2} (\hat{n} \cdot \vec{\nabla}\theta(\vec{x}_i))^2
\end{equation}
The interaction energy for this single bond is:
\begin{equation}
    -J \cos(\theta_i - \theta_j) \approx -J \left(1 - \frac{a^2}{2} (\hat{n} \cdot \vec{\nabla}\theta)^2\right) = \text{const} + \frac{Ja^2}{2} (\hat{n} \cdot \vec{\nabla}\theta)^2
\end{equation}
The constant term $-J$ simply shifts the total energy by a fixed amount and does not affect the physics of the fluctuations, so we can drop it.

\subsection{From a Discrete Sum to a Continuous Integral}
The total Hamiltonian is the sum over all bonds in the lattice. It is more convenient to rewrite this as a sum over sites. Since each bond connects two sites, we can sum over sites and divide by two to avoid double counting.
\begin{equation}
    H \approx \frac{1}{2} \sum_i \sum_{\hat{n}} \frac{Ja^2}{2} (\hat{n} \cdot \vec{\nabla}\theta_i)^2
\end{equation}
where the second sum is over all unit vectors $\hat{n}$ pointing to the neighbors of site $i$. For a hypercubic lattice, these are the vectors $\pm \hat{x}_1, \dots, \pm \hat{x}_d$. The sum over the gradient term gives:
\begin{align}
    \sum_{\hat{n}} (\hat{n} \cdot \vec{\nabla}\theta_i)^2 &= \sum_{k=1}^d \left[ (\hat{x}_k \cdot \vec{\nabla}\theta)^2 + (-\hat{x}_k \cdot \vec{\nabla}\theta)^2 \right] \\
    &= \sum_{k=1}^d \left[ (\partial_k \theta)^2 + (\partial_k \theta)^2 \right] = 2 \sum_{k=1}^d (\partial_k \theta)^2 = 2 (\vec{\nabla}\theta)^2
\end{align}
Plugging this back into the Hamiltonian:
\begin{equation}
    H \approx \frac{1}{2} \sum_i \left( \frac{Ja^2}{2} \cdot 2(\vec{\nabla}\theta_i)^2 \right) = \sum_i \frac{Ja^2}{2} (\vec{\nabla}\theta_i)^2
\end{equation}
This is our effective Hamiltonian, still written as a discrete sum over lattice sites. The final step is to take the continuum limit. In this limit, the sum over discrete sites can be replaced by a continuous integral over space, provided we remember that the density of sites is $1/a^d$:
\begin{equation}
    \sum_i \quad \longrightarrow \quad \int \frac{d^d x}{a^d}
\end{equation}
Our Hamiltonian becomes the \textbf{free energy functional}, $F[\theta(\vec{x})]$:
\begin{equation}
    F[\theta(\vec{x})] = \int \frac{d^d x}{a^d} \left( \frac{Ja^2}{2} (\vec{\nabla}\theta(\vec{x}))^2 \right) = \int d^d x \frac{J a^{2-d}}{2} (\vec{\nabla}\theta(\vec{x}))^2
\end{equation}
By defining a new constant $K = J a^{2-d}$, which absorbs all the microscopic details, we arrive at our final, beautiful result:
\begin{equation}
    F[\theta(\vec{x})] = \int d^d x \frac{K}{2} (\vec{\nabla}\theta(\vec{x}))^2
    \label{eq:free_energy_final}
\end{equation}
This is the action for a \textbf{free, massless scalar field theory}. It tells us that the "energy" or "cost" of a particular field configuration is proportional to how much it varies in space. Smooth, constant fields have zero energy, while rapidly fluctuating fields have a high energy cost. This is the essence of a field theory.

\newpage
\section{The Path Integral and Core Principles}
We have derived the effective energy functional for our system. How do we use it to calculate physical quantities? We must generalize the partition function to the continuum.

\subsection{The Functional (Path) Integral}
In discrete statistical mechanics, the partition function is a sum over all possible configurations of the system:
\begin{equation}
    Z = \sum_{\{\theta_i\}} e^{-\beta H[\{\theta_i\}]} = \left( \prod_i \int_0^{2\pi} \frac{d\theta_i}{2\pi} \right) e^{-\beta H[\{\theta_i\}]}
\end{equation}
In the continuum limit, a "configuration" is no longer a set of discrete numbers $\{\theta_i\}$, but a continuous function (a field) $\theta(\vec{x})$ defined over all of space. The sum over all configurations becomes an integral over the infinite-dimensional space of all possible field configurations. This is called a \textbf{functional integral} or, in a quantum context, a \textbf{path integral}:
\begin{equation}
    Z = \int \funcint{\theta(\vec{x})} e^{-\beta F[\theta(\vec{x})]}
\end{equation}
Here, $\int \funcint{\theta(\vec{x})}$ is the symbolic notation for this extraordinary object. Using our derived free energy (\ref{eq:free_energy_final}), the partition function for the effective field theory is:
\begin{equation}
    Z = \int \funcint{\theta(\vec{x})} \exp\left( -\frac{\beta K}{2} \int d^d x \, (\vec{\nabla}\theta(\vec{x}))^2 \right)
\end{equation}
This is the central object of study in statistical field theory. While technically challenging to evaluate, it contains all the information about the long-distance physics of our original spin model. The probability of observing a particular field configuration $\theta(\vec{x})$ is given by $p[\theta(\vec{x})] = \frac{1}{Z} e^{-\beta F[\theta(\vec{x})]}$.

\subsection{The Four Pillars of Field Theory}
Our journey from a discrete model to a continuum field theory has revealed the foundational principles that underpin all such theories, from condensed matter to the Standard Model.

\begin{enumerate}
    \item \textbf{Fluctuations:} The system exhibits fluctuations at all scales, especially near criticality. The path integral is precisely the mathematical tool designed to sum over all these possible fluctuations to compute physical observables. It embodies the principle of exploring all possibilities, weighted by their probability.

    \item \textbf{Locality:} The interactions are \textbf{local}. The energy cost at a point $\vec{x}$ depends only on the field and its derivatives \textit{at that same point}. This is manifest in our action, which is an integral of a Lagrangian density, $\mathcal{L} = \frac{K}{2} (\vec{\nabla}\theta)^2$. There are no terms that directly couple the field at $\vec{x}$ to the field at a distant point $\vec{y}$. All interactions are mediated through the local structure of spacetime.

    \item \textbf{Symmetry:} Symmetries of the microscopic system often survive in the continuum limit and have profound consequences. Our original Hamiltonian was invariant under a global rotation of all spins, $\theta_i \to \theta_i + \alpha$. This global U(1) symmetry persists in our continuum action, which is invariant under a constant shift $\theta(\vec{x}) \to \theta(\vec{x}) + \alpha$, because the derivatives $\vec{\nabla}\theta$ are unchanged. In QFT, \textbf{Noether's theorem} states that every continuous symmetry of the action corresponds to a conserved quantity.

    \item \textbf{Universality (Renormalizability):} This is perhaps the deepest lesson. The long-distance, collective behavior of a system near a critical point is \textbf{universal}—it is insensitive to the microscopic details. Our final action (\ref{eq:free_energy_final}) has a single parameter, $K$, that bundles up all the messy details of the original model ($J$, $a$, the lattice structure). Different microscopic models (e.g., spins on a triangular lattice, or even a liquid-gas system) can flow to the exact same effective field theory at long distances. This is why field theory is so powerful: it allows us to make precise predictions about collective behavior without needing to know the ultimate "theory of everything" at the shortest scales.
\end{enumerate}

\newpage
\section{Conclusion and Outlook}

In this lecture, we have followed a complete logical path from a simple, classical many-body system to a continuous field theory.

\textbf{Our Journey:}
\begin{itemize}
    \item We started with a discrete model of "particles" (classical spins on a lattice) with local interactions.
    \item We recognized that near a continuous phase transition, the correlation length diverges, making the discrete lattice structure irrelevant.
    \item We performed a systematic coarse-graining and continuum limit, assuming the system's configurations are slowly varying at long distances.
    \item This process led us to an effective description in terms of a continuous field, $\theta(\vec{x})$, whose dynamics are governed by a free energy functional, or action.
    \item The statistical mechanics of this field are described by a functional integral over all possible field configurations.
\end{itemize}

This demonstrates that field theory is not an abstract construction imposed upon nature, but rather the \textit{inevitable emergent description} of any many-body system whose correlations are much larger than its microscopic scales. The very same principles apply when moving from statistical field theory to quantum field theory. The main difference is that we will be dealing with fields in $(d+1)$-dimensional spacetime instead of $d$-dimensional space, and the path integral will involve a quantum action $S$ instead of a classical energy $\beta F$.

\subsection{Questions for Further Thought}
This new perspective invites us to reconsider what we mean by "fundamental".
\begin{enumerate}
    \item \textbf{Quasi-particles vs. Fundamental Particles:} In condensed matter, collective excitations of a field (like a spin wave in our magnet) are called quasi-particles. How are they different from "fundamental" particles like electrons?
    \item \textbf{The Primacy of Fields:} Does the existence of quasi-particles and the success of our emergent field theory suggest that fields are a more fundamental description of nature than particles?
    \item \textbf{Defining a Particle:} If we start with a field theory as our fundamental axiom, how do we define what a "particle" is?
    \item \textbf{Theories Without Particles:} Are there interacting field theories that have no well-defined particle-like excitations at all? (The answer is yes!)
\end{enumerate}

These questions will be at the heart of our exploration as we move forward to quantize these fields and build the framework of modern physics.

\end{document}